\color{uniblue}\section[Выбор уставок защиты от потери охлаждения]{\sloppy Выбор уставок защиты от потери охлаждения} 
\color{black}

\begin{enumerate}[label=\arabic{section}.\arabic*, labelsep=4pt, leftmargin=0pt, itemindent=57pt]

\item
Уставку срабатывания ИО тока обратной последовательности $I_2$ рекомендуется принимать равной 25\% от номинального тока. Указанное значение обеспечивает надежную отстройку от допустимой максимальной несимметрии питающей сети и погрешности ТТ при протекании через трансформатор малых значений тока нагрузки.
\item\label{zop:punkt}
Уставка срабатывания ИО по отношению тока обратной последовательности к току прямой последовательности $K_\textsl{несим}$ может быть принята в диапазоне от 0,3 до 0,4 для срабатывания ЗОП как при обрыве в первичной сети, так и при во вторичных цепях.
\item
Выбор режима работы ЗОП осуществляется с помощью программного переключателя <<Реж\_ЗОП>>, для выбора доступны следующие возможные состояния:
\begin{itemize}
\item по току обратной последовательности (ИО I2);
\item по отношению тока обратной последовательности к току прямой последовательности (ИО I2/I1);
\item по обоим ИО (ИО I2 + ИО I2/I1).
\end{itemize}

\item
Выдержка времени выбирается на ступень селективности больше времени срабатывания самой медленно действующей защиты трансформатора:
\begin{equation}
    \label{eq:zop}
    t_{\textsl{с.з.ЗОП}} = t_{\textsl{с.з.МТЗ}}+{\Delta t},
\end{equation}
\begin{eqexpl}[15mm]
    \item{$t_{\textsl{с.з.МТЗ}}$} время срабатывания самой медленно действующей защиты питающей сети, с;
    \item{$\Delta t$} ступень селективности (принимается равной 0,3-0,5), с.
\end{eqexpl}

\end{enumerate}

